\chapter*{Pengantar}

\begin{center}
\textarabic{﷽}
\end{center}

Kepada-Nya aku bersandar, kepada-Nya rasa \textit{tsiqah} percayaku berlabuh, dan kepadanya pula aku memohon pertolongan atas seluruh urusan dunia dan agama.

Segala puji bagi Allah yang telah menunjukkan kita untuk mengenal yang halal dan haram, dan mewajibkan atas kita untuk taat kepada Nabi kita Muhammad, pemimpin seluruh makhluk. Dia juga memerintahkan kita untuk mengikuti syariat-Nya yang agung dan memahami hukum-hukum, serta memperbolehkan kita memandang Wajah-Nya Yang Maha Mulia di Darussalam. Dia Allah membenci bagi kita kekufuran, kefasikan, dan kemaksiatan, serta mengharamkan atas kita perbuatan dosa dan kesewenang-wenangan.

Aku bersaksi bahwa tidak ada \textit{illah} yang berhak diibadahi selain Allah semata, tiada sekutu bagi-Nya. Pemilik keagungan, kemuliaan, keperkasaan, dan kekuasaan. Aku bersaksi bahwa Muhammad adalah hamba dan utusan-Nya, sang penyeru kepada pendalaman fikih agama. Semoga shalawat dan salam Allah tercurahkan kepadanya, keluarganya, dan para sahabatnya sekalian, \textit{amma ba’du}.

Risalah ini adalah buku ringkas dalam memilih pendapat-pendapat paling kuat dari para ulama berkaitan dengan sebagian permasalahan-permasalahan thaharah dari kitab matan \textbf{Safinatun Najah} dan \textbf{Nailur Raja’ Bisyarhi Safinatun Naja}. Buku ini dikhususkan bagi santriwati setingkat SMP. Saya menamai buku ini dengan judul \textbf{Syarhu Safinatin Najati Fi Ba’dhi Ahkami Thaharati Li Thalibat Mustawa Tsanawi}.

Saya menyusunnya sebagai pendekatan pengantar bagi pemula dengan mencukupkan pada isyarat terhadap dalil, tanpa menyelami rincian, bantahan serta uraian yang panjang. Masalah-masalah fiqih yang dibahas merupakan pendapat-pendapat dari madzhab Syafi’i \textit{rahimahullah}, yang diambil dari pendapat beliau, pilihan ulama pengikut beliau dan ulama madzhab yang masyhur. Dalam mengambil pendapat, saya tidak terikat dengan pendapat masyhur dari ulama \textit{muta’akhirin} (kontemporer), akan tetapi pegangan dalam hal ini mengacu pada pendapat yang memiliki dalil, berdasarkan penelitian Syeikh Abu Qatadah Al-filistini \textit{hafidzahullah} dalam kitabnya \textbf{Mukhtashar Abi Qatadah Fi Fiqhi Asy-Syafi'iyah As-Sadah}.

Syeikh Syeikh Abu Qatadah Al-filistini \textit{hafidzahullah} mengatakan: 

"Berkomitmen pada suatu madzhab fikih bukan berarti menolak dalil dan tarjih (memilih pendapat yang paling kuat). Madzhab-madzhab ulama hakikatnya diletakkan untuk mendekatkan terhadap dalil dan penjelasan terhadap pendapat imam madzhab, para pengikutnya, dan \textit{ashabul wujuh } (ulama yang mengembangkan madzhab) bukan supaya taklid kecuali bagi orang awam yang keadaannya tidak memungkinkan untuk selain itu.Kemudian tidak ada suatu madzhab kecuali di dalamnya terdapat ruang yang luas untuk berijtihad dan memilih pendapat. Inilah tujuan saya setelah mempermudah masalah-masalah fikih dan meringkasnya, dan Allahlah yang memberikan taufik.

Bagi siapa yang membaca dan menelaah madzhab Syafi’i akan mengetahui keluasannya, dan sesungguhnya para ulama madzhab tidaklah mengambil pendapat Imam Syafi’i kecuali memiliki dasar dalil untuk berijtihad dan memilih pendapat bukan untuk taqlid yang justru dibenci oleh imam Syafi’i itu sendiri yang telah melarang hal tersebut seperti yang disebutkan oleh Immam  Al-Muzani \textit{rahimahullah} dalam Mukhtasarnya. Maka barang siapa yang bertaqlid kepada Syafi’i tanpa penelitian, ilmu, dan pengetahuan akan dalilnya sesungguhnya dia telah membatalkan tujuan utama Imam Syafi’i dalam ilmu ini. Hal ini sepantasnya harus dipahami dalam perkara ini. Tidak diperbolehkan bagi seseorang untuk menolak pendapat orang lain dengan alasan bahwa pendapat yang ia pilih tersebut menyelisihi pendapat yang masyhur dalam madzhab. Padahal pendapat yang masyhur itu sebenarnya hanyalah ijtihad satu atau lebih para ulama. Sedangkan dia hanya sedang memilih suatu pendapat dari para ulama -Selesai perkataan.

Allah tempat meminta, agar menjadikan karya ini ikhlas semata-mata hanya karena mengharapkan wajah Allah yang mulia, dan menjadikan kitab ini tersebar luas seperti peredaran matahari di penjuru ufuk, serta memberikan manfaat pada penulisnya dan para penghafalnya, dan orang yang menyibukkan diri dengannya. Tiada taufik bagiku kecuali dengan pertolongan dari Allah, hanya kepada-Nya aku bertawakal dan bertaubat. Dialah Allah sebaik-baik pelindung dan penolong. Maha suci Rabbmu, Rabb pemilik kemuliaan dari apa yang mereka sifatkan, dan salam atas seluruh rasul, \textit{walhamdulillahi rabbil 'alamin}.

\vspace{1cm} % Memberi jarak vertikal antara teks dan blok penulis

\begin{flushleft}
Ditulis oleh, hamba yang paling fakir pada Allah \textit{ta'ala}\\
\textbf{Zen Ibrahim} semoga Allah mengampuninya dan kedua orang tuanya.\\
\textbf{Markaz Mulazamah Mutiara Muslimah}\\
Karanganyar - Jawa Tengah\\
Hari Kamis 14 Syawal 1447 H bertepatan dengan 2 April 2026 M
\end{flushleft}